\section{Una perspectiva unificada de los modelos basados en puntuaciones y DDPM}

\subsection{La convergencia de dos vías de investigación}

Históricamente, DDPM y los modelos basados en puntuaciones (score-based models) fueron dos vías de desarrollo independiente:

\begin{knowledgebox}{El desarrollo de ambas vías}
\begin{itemize}
    \item \textbf{Vía DDPM}: Ho et al. (2020) derivaron el objetivo de entrenamiento basado en inferencia variacional y ELBO. Partiendo de las tradiciones de los modelos gráficos probabilísticos y los autoencoders variacionales, definieron cadenas de Markov hacia adelante/hacia atrás y entrenaron el proceso inverso maximizando el ELBO.
    \item \textbf{Vía basada en puntuaciones}: Song \& Ermon (2019, 2020) se basaron en score matching y dinámicas de Langevin. Partiendo de la función score en estadística, entrenaron un estimador de score mediante denoising score matching y realizaron muestreo usando dinámicas de Langevin con recocido (annealed).
\end{itemize}
Ambas vías finalmente se unificaron en el marco de SDE basado en puntuaciones de Song et al. (2021).
\end{knowledgebox}

\subsection{Resumen matemático de la relación de equivalencia}

Ya hemos establecido progresivamente todas las relaciones de equivalencia en las secciones anteriores; aquí hacemos un resumen sistemático:

\begin{enumerate}
    \item \textbf{Equivalencia del objetivo de entrenamiento}: La pérdida de predicción de $\epsilon$ en DDPM es equivalente a la pérdida de Denoising Score Matching (con un factor de escala diferente).
    \item \textbf{Equivalencia de la salida de la red}: $\epsilon_\theta(x_t, t) = -\sqrt{1-\bar{\alpha}_t}\, s_\theta(x_t, t)$.
    \item \textbf{Equivalencia del proceso de muestreo}: El muestreo inverso en DDPM es equivalente a las dinámicas de Langevin con recocido (en su forma discretizada).
    \item \textbf{Conexión Tweedie}: $\mathbb{E}[x_0 \mid x_t]$ puede expresarse tanto mediante la función score como mediante un predictor de ruido.
\end{enumerate}

\begin{importantbox}{Significado práctico de la perspectiva unificada}
Comprender esta equivalencia no es solo una cuestión de elegancia académica, sino que tiene implicaciones prácticas importantes:
\begin{itemize}
    \item Se puede utilizar la teoría basada en puntuaciones para analizar las propiedades de DDPM (como la convergencia y la calidad del muestreo).
    \item Se pueden emplear las técnicas prácticas de DDPM para mejorar el muestreo basado en puntuaciones (por ejemplo, en el diseño del schedule).
    \item La perspectiva de la función score proporciona una base teórica natural para la generación condicional (guía por clasificador, guía libre de clasificador).
    \item El marco de SDE de tiempo continuo de Song et al. generaliza los pasos discretos a procesos continuos, abriendo el camino para diseños de muestreadores más flexibles (como DDIM, DPM-Solver, entre otros).
\end{itemize}
\end{importantbox}

\subsection{Resumen de la sección}

Aunque DDPM y los modelos basados en puntuaciones surgieron de diferentes puntos de partida y motivaciones, son matemáticamente completamente equivalentes. DDPM ofrece una perspectiva de modelos gráficos probabilísticos e inferencia variacional, mientras que los modelos basados en puntuaciones aportan la intuición física de la función score y las dinámicas de Langevin. Comprender ambas perspectivas es crucial para dominar profundamente la teoría de los modelos de difusión.


